\section{Fréquence d'une valeur}

\begin{definitionbox}
    La \bb{fréquence d'une valeur} est le quotient de l'effectif de cette
    valeur par l'effectif total.
    \begin{center}
        $\text{Fréquence} = \dfrac{\text{Effectif}}{\text{Effectif total}}$
    \end{center}
\end{definitionbox}

\begin{exempleboxsuite}
    Pour \bb{calculer la fréquence} des élèves venant à pied au collège on
    connait :

\begin{itemize}[label=\textbullet]
    \item l'\bb{effectif} des élèves venant à pied : \answer{125}
    \item l' \bb{effectif total} : \answer{500}
\end{itemize}

    On en déduit la fréquence des élèves venant à pied : $f$ = 
    \answer{$\dfrac{125}{500} = \dfrac{1}{4} = 0.20 = 20\%$}
\end{exempleboxsuite}

\begin{remarquebox}
    Une fréquence peut se donner
    \begin{itemize}[label=\textbullet]
        \item sous forme d'une fraction simplifiée : $f$ = \answer{$\dfrac{1}{4}$}
        \item en écriture décimale : $f$ = \answer{$0.25$}
        \item en pourcentage : $f$ = \answer{$25\%$}
    \end{itemize}

    \vspace{20pt}

    On peut également présenter les résultats dans un \bb{tableau des
    fréquences} :

    \begin{center}
    \begin{tabular}{lccccc}
       \bb{Transport} & \bb{2 roues} & \bb{Bus} & \bb{À pied} & \bb{Parents} & \bb{Total}\\
       \midrule
       \bb{Effectif} & \answer{50} & \answer{250} & \answer{125} & \answer{75} & \answer{500} \\
       \bb{Fréquence} & \answer{$\frac{50}{500}=0.1$} &
        \answer{$\frac{250}{500}= 0.5$}
        & \answer{$\frac{125}{500}$} & \answer{$\frac{75}{500}=0.15$} &
        \answer{$500$} \\
       \bb{Fréquences en \%} & \answer{$10\%$} & \answer{$50\%$} &
        \answer{$25\%$} & \answer{$15\%$} & \answer{$100\%$} \\
    \end{tabular}
    \end{center}
\end{remarquebox}

\begin{remarquebox}
    Une fréquence est toujours comprise entre \answer{$0$} et \answer{$1$} \\

    La somme des fréquences est égale à \answer{$1$} ou (ou à \answer{$100\%$})

\end{remarquebox}

